The math for this project will be split into three parts:
\begin{enumerate}
	\item Calculating the survivor fraction
	\item coordinate randomization
	\item mass distribution 
\end{enumerate}
\subsection{calculating the survivor fraction}
The survivor fraction is how we measure whether the Uranium has reached the critical mass. It is calculated by dividing the previous "generation" of fission by the current "generation" of fission. 

If the survivor fraction ($f$) is lower than 1, we can conclude that we are below the critical mass, meaning that the reaction will die out.

If the survivor fraction is higher than 1, we can conclude that the reaction has gone supercritical.

By changing the mass (user input), we can try to find the point at which the survivor fraction is 1. To make sure this is accurate, we will allow the user to insert an amount of points ($N$) that need to be simulated. we will then simulate those $N$ random points inside the volume. this will be our first generation. We will then simulate fission for each of those $N$ points. If we check whether the neutrons are inside the volume, we can calculate the survivor fraction with the formula:
\begin{equation}
	f=\frac{N_{\text{inside volume}}}{N_{\text{gen 1}}}
\end{equation}
To be able to check whether the neutrons are inside the volume, we first need to build a randomizer for the movement of the neutrons inside the simulations.
\subsection{coordinate randomization}
Using a randomizer to generate random uniformly distributed points, we can calculate the path of the neutrons. If the neutrons have starting coordinates ($x_0$, $y_0$, $z_0$), we can describe their endpoints ($x_1$, $y_1$, $z_1$) using the equations:
\begin{displaymath}
	\begin{cases}
		x_1= x_0+d\;\sin\theta\;\cos\phi \\
		y_1=y_0+d\;\sin\theta\;\sin\phi \\
		z_1=z_0+d\;\cos\theta
	\end{cases}
\end{displaymath}